% !TeX root = ../papers/nc-p3.tex
\section{二つの三次系による楕円四次成分の排除}
\label{sec:rank-one-elliptic-quartic-exclusion}
\begin{theorem}\label{thm:rank-one-elliptic-quartic-exclusion}
四辺有限の標準有限Frobenius窓で、入力は被約、反対平方根は一般可逆で全二次式を成分に持つとする。
階数一軌跡の曲線成分は直線と平面二次曲線だけであり、その被約総次数は5以下である。
\end{theorem}
\begin{proof}[二つの三次系の積を比較する]
残る整四次成分 $E$ は前定理で滑らかな楕円正規四次、二埋込みの線束は $L,L'=L\eta$、
$\eta$ は非自明な二捩れである。一般点 $a\in E$ と $\mathcal O(b-a)=\eta$ となる $b$ を選ぶ。
次数三の二束を
\[
 A=L(-a)\simeq L'(-b),\qquad B=L(-b)\simeq L'(-a)
\]
とする。$A,B$ はそれぞれ三切断であり、標準切断による乗法でどちらも両座標四空間に入る。
$A$ の二つの包含を $i,i'$、$B$ のものを $j,j'$ とすれば、全ての $s\in H^0(A),u\in H^0(B)$ で
$i(s)j(u)$ と $i'(s)j'(u)$ は同じ $a+b$ の切断を掛けたものとなる。
二乗束の同一視と包含の定数を合わせて、後者を反対二次座標同型で入力に戻したものと前者の差は
$E$ を切る二つの二次式 $q_1,q_2$ の包に属する。
従ってこれらの差は $s,u$ の基底で三行三列の一次行列鉛筆 $P(q_1,q_2)$ になる。対称性は仮定しない。

各曲線外印の入力点を $x$、出力点を $y$ とし、$i,i'$ の三切断の値を $\alpha_x,\alpha'_y$、
$j,j'$ の値を $\beta_x,\beta'_y$ と書く。点ごとの自明化による非零定数を $c$ として
\[
          P(q(x))=\alpha_x\beta_x^t-c\,\alpha'_y(\beta'_y)^t .
\]
特に階数は2以下である。曲線外印は少なくとも八つ、同一の鉛筆値での上界は一群を除いて二つ、
その一群でも三つなので、値は少なくとも四種類ある。従って三次式 $\det P$ は恒等零である。
\end{proof}
\begin{proof}[左右の原始核列の次数]
一般の $a$ では、各指定印で $\alpha_x\wedge\alpha'_y$ と $\beta_x\wedge\beta'_y$ は両方零か
両方非零である。実際 $a$ を $b=a+\eta$ に移す対合は $A,B$ とその両包含を定数倍まで交換する。
固定印での零条件は代数的である。例えば対角を除いた族 $L(-a)$ の切断を底へ押し出すと階数三束になり、
二つの包含は底の線束を局所自明化して書ける。一方の条件が底全体なら他方も全体で、
そうでなければ双方とも有限集合である。有限個の印について同時に一般点を取れる。

すべての印での $\alpha_x\wedge\alpha'_y$ は三次元を張らなければならない。
さもなくば二つの三切断空間における非零の交代テンソルが全印で零値となり、
両包含が単射なので十六の混合積の独立性に反する。曲線上の印では二つの評価は同じ三次系の評価の定数倍であり
（片方が零の場合も含む）、外積は零である。同じ理由で $\beta_x\wedge\beta'_y$ も曲線外で三次元を張る。
非零外積のある印では両側とも独立な二列なので $P(q(x))$ の階数はちょうど2である。
その左核は $\alpha_x,\alpha'_y$ の零化直線、右核は $\beta_x,\beta'_y$ の零化直線である。
従って一般階数は2であり、二種類の核列の値はそれぞれ三次元全体を張る。

ところが $\operatorname{adj}P$ は非零の階数一の二項二次式行列で、原始な二項多項式列 $g,h$ と
スカラー多項式 $d$ により $d g h^t$ と書ける。一意分解により分母がないため
$\deg g+\deg h\le2$ である。階数2の値ではこれらが左右の核列を与える。
各々の値が三次元を張ればそれぞれ次数が2以上必要であり、次数和に反する。
これで楕円四次も排除された。次数三以上の他の成分は既に除かれたので結論を得る。
\end{proof}
